\documentclass[11pt,a4paper]{article}

% -----------------------------------------------------------------------------
% Paquetes: todos forman parte de una instalacion estandar de LaTeX.
% -----------------------------------------------------------------------------
\usepackage[utf8]{inputenc}
\usepackage[T1]{fontenc}
\usepackage{enumitem}
\usepackage{amsmath}
\usepackage{amssymb}
\usepackage{booktabs}
% Fecha en espanol sin depender de babel ni de spanish.ldf.
\newcommand{\FechaEnEspanol}{%
  \number\day\space de\space
  \ifcase\month
    \or enero\or febrero\or marzo\or abril\or mayo\or junio\or
    julio\or agosto\or septiembre\or octubre\or noviembre\or diciembre
  \fi\space de\space\number\year
}
\usepackage[a4paper,margin=2.5cm,headheight=42pt]{geometry}
\usepackage{graphicx}
\usepackage{xcolor}
\usepackage{fancyhdr}

% lastpage no es necesario para compilar, pero permite mostrar "pagina de N".
\newif\ifHayUltimaPagina
\IfFileExists{lastpage.sty}{
  \usepackage{lastpage}
  \HayUltimaPaginatrue
}{
  \HayUltimaPaginafalse
}

% -----------------------------------------------------------------------------
% CONFIGURACION DEL TRABAJO: editar solamente este bloque para reutilizarlo.
% -----------------------------------------------------------------------------
\newcommand{\NombreMateria}{Redes de Computadoras}
\newcommand{\TipoTrabajo}{Trabajo Práctico}
\newcommand{\NumeroTrabajo}{3}
% Cambiar por "Teórico" o "Práctico" según corresponda.
\newcommand{\TipoTP}{Teórico}
\newcommand{\NombreTrabajo}{Visión general de las redes de área local}
\newcommand{\NombreGrupo}{NetRunners}
\newcommand{\Docente}{FACUNDO NICOLAS OLIVA CUNEO}
\newcommand{\LogoArchivo}{logo.png}
\newcommand{\LogoRuta}{\LogoArchivo}

% Una linea por integrante. Se pueden agregar o quitar filas.
\newcommand{\Integrantes}{%
  \begin{tabular}{ll}
    Bastida, Lucas Ramiro & DNI 39444511 \\
    Contessi, Ruben Andres & DNI 33315235 \\
    Garay, Ignacio Hernan & DNI 44191925 \\
    Giraudo, Juan Pablo & DNI 34970089 \\
    Peñaloza, Gonzalo Adrian & DNI 40441355 \\
    Quinteros del Castillo, Tomás Agustín & DNI 46169739 \\
    Vasconsellos Blason, Benjamin Alejandro & DNI 45642879 \\
  \end{tabular}%
}

% -----------------------------------------------------------------------------
% Elementos comunes de caratula y encabezado.
% -----------------------------------------------------------------------------
\newcommand{\InsertarLogo}[1][0.24\textwidth]{%
  \ifx\LogoRuta\empty
    \fbox{\parbox[c][1cm][c]{#1}{\centering Logo institucional\\\small (colocar archivo)}}%
  \else
    \includegraphics[width=#1,keepaspectratio]{\LogoRuta}%
  \fi
}

\newcommand{\TituloCorto}{TP \NumeroTrabajo (\TipoTP): \NombreTrabajo}

\pagestyle{fancy}
\fancyhf{}
\fancyhead[L]{\InsertarLogo[1.4cm]}
\fancyhead[C]{\small\textbf{\TituloCorto}}
\fancyhead[R]{\small\NombreGrupo}
\renewcommand{\headrulewidth}{0.4pt}
\fancyfoot[C]{\small Pagina \thepage\ifHayUltimaPagina\ de \pageref{LastPage}\fi}
\renewcommand{\footrulewidth}{0pt}

\begin{document}

\begin{titlepage}
  \thispagestyle{empty}
  \begin{center}
    \InsertarLogo[0.32\textwidth]

    \vfill
    {\Large\textsc{\NombreMateria}\par}
    \vspace{1.5cm}
    {\Huge\bfseries \TipoTrabajo\par}
    \vspace{0.35cm}
    {\LARGE N.\textsuperscript{o} \NumeroTrabajo (\TipoTP)\par}
    \vspace{0.8cm}
    {\Large\NombreTrabajo\par}

    \vfill
    \begin{tabular}{rl}
      \textbf{Grupo:}   & \NombreGrupo \\
      \textbf{Docente:} & \Docente
    \end{tabular}

    \vspace{1cm}
    \begin{center}
      \textbf{Integrantes}\\[0.2cm]
      \Integrantes
    \end{center}

    \vfill
    {\large \FechaEnEspanol\par}
  \end{center}
\end{titlepage}

\section*{Cuestiones de repaso}

\begin{enumerate}[label=15.\arabic*), itemsep=0.3cm]

  \item
  \item
  \item
  \item
  \item
  \item
  \item
  \item
  \textbf{Enumere y describa brevemente los modos de operación proporcionados por el protocolo LLC.}

        El protocolo LLC (Control de Enlace Lógico) define tres modos de operación para gestionar el intercambio de datos entre los puntos de acceso al servicio (SAP) de la capa de enlace:

        \begin{enumerate}[label=\arabic*., itemsep=0.2cm]
          \item \textbf{Tipo 1: servicio no orientado a conexión sin confirmación.} Es un servicio de datagramas que no incluye mecanismos de control de flujo ni de control de errores, por lo que la recepción de los datos no está garantizada a este nivel. Los datos se transmiten mediante PDU de información no numerada (UI).

          \item \textbf{Tipo 2: servicio orientado a conexión.} Establece una conexión lógica entre los SAP de las estaciones antes de transferir datos. Utiliza el modo balanceado asíncrono (ABM) de HDLC y proporciona control de flujo, control de errores y numeración secuencial de tramas.

          \item \textbf{Tipo 3: servicio no orientado a conexión con confirmación.} No requiere establecer previamente una conexión lógica, pero cada PDU de datos enviada es confirmada individualmente por el receptor. Utiliza PDU de información no orientada a conexión con confirmación (AC) y un número de secuencia de 1 bit para el control de entrega.
        \end{enumerate}
        
  \item
  \item
  \item
  \item 
  \textbf{¿Qué diferencias existen entre un concentrador y un conmutador de capa 2?}

        Aunque ambos dispositivos se utilizan como nodo central en topologías en estrella, presentan diferencias en su funcionamiento y rendimiento:

        \begin{itemize}
          \item \textbf{Funcionamiento y tratamiento de la señal:} El concentrador (\textit{hub}) opera en la capa física como un repetidor multipuerto: retransmite la señal recibida por un puerto hacia todos los demás, sin analizar la dirección de destino. El conmutador de capa 2 (\textit{switch}) opera en la capa de enlace de datos: analiza la dirección MAC de destino de la trama y, si conoce el puerto asociado, la reenvía únicamente por ese puerto.

          \item \textbf{Capacidad del medio y colisiones:} En un concentrador, todas las estaciones comparten el ancho de banda y un mismo dominio de colisión; solo una puede transmitir a la vez sin provocar colisiones. En un conmutador, cada puerto constituye un dominio de colisión separado y dispone de capacidad dedicada. Esto permite varias comunicaciones simultáneas entre puertos distintos y, en enlaces \textit{full-duplex}, elimina las colisiones.

          \item \textbf{Procesamiento y modos de transmisión:} El concentrador no almacena ni procesa tramas individualmente. El conmutador realiza el reenvío mediante hardware especializado y puede gestionar búferes de tramas. Según el equipo, puede funcionar mediante almacenamiento y envío (\textit{store-and-forward}) o conmutación rápida (\textit{cut-through}), además de admitir enlaces \textit{full-duplex}.
        \end{itemize}
        
  \item

\end{enumerate}

\end{document}
